% =========================================================
%  General Formularium – Physics
%  Chapter 4: Statics & Fluid Mechanics
%  Converted from "General Formularium (ENG) - V2.docx"
%  Reviewed and corrected on 21 Oct 2026.
% =========================================================
\documentclass[11pt]{article}
\usepackage[a4paper,margin=2cm]{geometry}
\usepackage{amsmath,amssymb}
\usepackage[dvipsnames]{xcolor}
\setlength{\parindent}{0pt}
\newcommand{\sect}[2]{\section*{\color{#1}#2}}
\newcommand{\sub}[1]{\par\medskip\textbf{#1}\par}
\newcommand{\expl}[1]{{\color{gray}\itshape #1}\par}
\newcommand{\warn}[1]{{\color{BrickRed}\textbf{!} #1}}
\definecolor{cStF}{HTML}{1E8449}

\begin{document}

\sect{cStF}{Statics \& Fluid Mechanics}
\expl{Equilibrium, moments, distributed loads, and pressure in fluids and gases. Your Statics tab has more detail on statics. The unit is given in square brackets.}

% ---------- Statics ----------
\sub{Static equilibrium}
\[ \sum \vec F_i = \vec 0 \quad [\mathrm{N}] \qquad \sum \vec M = \vec 0 \quad [\mathrm{N\,m}] \]

\sub{Force equilibrium of a particle}
\[ \sum \vec F = \vec F_1 + \vec F_2 + \vec F_3 + \dots = \vec 0 \qquad [\mathrm{N}] \]

\sub{Moment of a force}
\[ \vec M_O = \vec r\times\vec F \qquad M_O = F\,d \qquad [\mathrm{N\,m}] \]

\sub{Moment about an axis}
\[ \vec M_l = M_l\,\hat u_l \qquad M_l = \hat u_l\cdot(\vec r\times\vec F) \qquad [\mathrm{N\,m}] \]

\sub{Resultant moment}
\[ \vec M_{R,O} = \sum (\vec r_i\times\vec F_i) + \sum \vec M_i \qquad [\mathrm{N\,m}] \]

\sub{Force component along an axis}
\[ \vec F_l = F_l\,\hat u_l \qquad F_l = \hat u_l\cdot\vec F = \|\vec F\|\cos\alpha \qquad [\mathrm{N}] \]

\sub{Perpendicular distance (lever arm) of the resultant}
\[ d = \frac{\left|\left(\vec M_R\right)_O\right|}{\left|\vec F_R\right|} \qquad [\mathrm{m}] \]

\sub{Line of action of parallel forces}
\[ x_R = \frac{M_O}{F_R} = \frac{\sum_i x_i F_i}{\sum_i F_i} \qquad [\mathrm{m}] \]

\sub{Resultant of a distributed load}
\[ F_R = \int_0^L w(x)\,dx \quad [\mathrm{N}] \qquad M_O = \int_0^L x\,w(x)\,dx \quad [\mathrm{N\,m}] \qquad x_R = \frac{\int_0^L x\,w(x)\,dx}{\int_0^L w(x)\,dx} \quad [\mathrm{m}] \]

\sub{Standard distributed loads}
\[ \text{uniform: } F_R = qL,\ \ x_R = \frac{L}{2} \qquad \text{triangular: } F_R = \frac{qL}{2},\ \ x_R = \frac{2L}{3} \]
\expl{For the triangle, $x_R$ is measured from the end where the load is zero.}

\sub{Shear force and bending moment}
\[ \frac{dV}{dx} = -w(x) \qquad \frac{dM}{dx} = V(x) \qquad V\ [\mathrm{N}],\ M\ [\mathrm{N\,m}] \]

\sub{Virtual work}
\[ \delta W = \sum F\,\delta r \qquad \text{equilibrium} \iff \delta W = 0 \qquad [\mathrm{J}] \]

\sub{Equilibrium and potential energy}
\[ \frac{dV}{dq} = 0 \]
\expl{$q$ = generalised coordinate (e.g. an angle). Stable if $\frac{d^2V}{dq^2} > 0$.}

\sub{Centroid, centre of mass, centre of gravity}
\[ \bar{\vec r} = \frac{\sum \tilde r_i\,w_i}{\sum w_i} \quad\text{(discrete)} \qquad \bar{\vec r} = \frac{\int \tilde r\,dw}{\int dw} \quad\text{(continuous)} \qquad [\mathrm{m}] \]
\expl{$w$ = weight factor: weight, mass, volume, area or length, depending on what you calculate.}

% ---------- Fluids ----------
\sub{Pressure}
\[ p = \frac{F}{A} \qquad p = \lim_{dA\to 0}\frac{dF}{dA} \qquad [\mathrm{Pa} = \mathrm{N/m^2}] \]

\sub{Density}
\[ \rho = \frac{m}{V} \qquad \rho = \lim_{dV\to 0}\frac{dm}{dV} \qquad [\mathrm{kg/m^3}] \]

\sub{Hydrostatic pressure}
\[ p = \rho\,g\,h \qquad [\mathrm{Pa}] \]

\sub{Gas pressure (U-tube manometer)}
\[ p_{\text{gas}} = p_{\text{atm}} + \rho\,g\,\Delta h \qquad [\mathrm{Pa}] \]

\sub{Buoyancy (Archimedes)}
\[ F_A = \rho_{\text{fluid}}\,g\,V_{\text{displaced}} \qquad [\mathrm{N}] \]

\sub{Bernoulli's principle}
\[ p + \tfrac{1}{2}\rho v^2 + \rho g h = \text{constant} \qquad [\mathrm{Pa}] \]

\sub{Ideal gas law}
\[ pV = nRT \qquad [\mathrm{Pa},\ \mathrm{m^3},\ \mathrm{mol},\ \mathrm{K}] \]

\sub{Combined gas law (any process)}
\[ \frac{p_1 V_1}{T_1} = \frac{p_2 V_2}{T_2} \]

\end{document}
